\section{Axiom Catalog}
\label{sec:axiom-catalog}

This appendix enumerates all axioms in the \MLIPsOntology{}, organized by
module. For each axiom we give the formal statement and a concrete example
illustrating its effect.

\subsection{Algorithm Module}

\subsubsection{Every algorithm has at least one hyperparameter}
\begin{align*}
  \MLIPAlgorithmC &\sqsub \existsR{hasHyperparameter}{Hyperparameter}
\end{align*}
\emph{Example.} The MACE algorithm declares hyperparameters such as cutoff
radius and number of layers. An $\MLIPAlgorithmC$ instance without any
$\hasHyperparameterR$ link would violate this axiom.

\subsubsection{Every algorithm supports at least one simulation type}
\begin{align*}
  \MLIPAlgorithmC &\sqsub \existsR{supportsSimulation}{SimulationType}
\end{align*}
\emph{Example.} MACE supports molecular dynamics and geometry optimization.
An algorithm that does not declare any supported simulation type is
incomplete.

\subsubsection{Every implementation specifies its library}
\begin{align*}
  \ImplementationC &\sqsub \existsR{implementedIn}{Library}
\end{align*}
\emph{Example.} ``MACE v0.3'' is an implementation in the MACE library.
An $\ImplementationC$ without an $\implementedInR$ link is invalid.

\subsubsection{Every training run executes exactly one algorithm}
\begin{align*}
  \TrainingRunC &\sqsub \exactR{1}{executes}{MLIPAlgorithm}
\end{align*}
\emph{Example.} A training run that produces a MACE model for Ti-Al
executes the MACE algorithm---not two algorithms simultaneously, and not
zero.

\subsubsection{Every training run uses exactly one training dataset}
\begin{align*}
  \TrainingRunC &\sqsub \exactR{1}{runsOn}{TrainingDataset}
\end{align*}
\emph{Example.} The training run uses the Ti-Al DFT dataset with 15{,}000
configurations. If a run combines multiple datasets, they should be
represented as a single merged $\TrainingDatasetC$.

\subsubsection{Every training run produces exactly one trained model}
\begin{align*}
  \TrainingRunC &\sqsub \exactR{1}{produces}{TrainedModel}
\end{align*}
\emph{Example.} The training run produces a single MACE model for Ti-Al.
If the same run produces checkpoints, only the final model is recorded.

\subsubsection{Every trained model was produced by some training run}
\begin{align*}
  \TrainedModelC &\sqsub \existsR{produces^{-}}{TrainingRun}
\end{align*}
\emph{Example.} A $\TrainedModelC$ instance ``MACE-TiAl-v1'' must be
linked to a $\TrainingRunC$ that produced it. A model cannot exist without
provenance about how it was trained.

\subsubsection{Every hyperparameter belongs to some algorithm}
\begin{align*}
  \HyperparameterC &\sqsub \existsR{hasHyperparameter^{-}}{MLIPAlgorithm}
\end{align*}
\emph{Example.} The hyperparameter ``cutoff radius'' is defined for the
MACE algorithm. A $\HyperparameterC$ instance that is not linked to any
algorithm via $\hasHyperparameterR^{-}$ is an orphan and violates this
axiom.

\subsubsection{Every implementation belongs to some algorithm}
\begin{align*}
  \ImplementationC &\sqsub \existsR{hasImplementation^{-}}{MLIPAlgorithm}
\end{align*}
\emph{Example.} ``MACE v0.3'' is an implementation of the MACE algorithm.
An $\ImplementationC$ that is not linked to any $\MLIPAlgorithmC$ via
$\hasImplementationR^{-}$ is invalid.

\subsubsection{Shortcut: trainedWith is a property chain}
\begin{align*}
  \trainedWithR &\equiv \producesR^{-} \circ \executesR
\end{align*}
\emph{Example.} If a $\TrainingRunC$ produces model $m$ and executes
algorithm $a$, then $m$ $\trainedWithR$ $a$ is inferred. This allows
querying ``which algorithm was used to train this model?'' without
navigating through the run.

\subsubsection{Shortcut: trainedOn is a property chain}
\begin{align*}
  \trainedOnR &\equiv \producesR^{-} \circ \runsOnR
\end{align*}
\emph{Example.} If a $\TrainingRunC$ produces model $m$ and runs on
dataset $d$, then $m$ $\trainedOnR$ $d$ is inferred. This allows
querying ``which dataset was used to train this model?'' directly.

\subsubsection{Domain axioms for algorithm-specific roles}
\begin{align*}
  \existsR{hasHyperparameter}{\top} &\sqsub \MLIPAlgorithmC \\
  \existsR{hasImplementation}{\top} &\sqsub \MLIPAlgorithmC \\
  \existsR{supportsSimulation}{\top} &\sqsub \MLIPAlgorithmC
\end{align*}
\emph{Example.} If an entity has a $\hasHyperparameterR$ link, it must be
an $\MLIPAlgorithmC$. This prevents, e.g., a $\TrainingDatasetC$ from
accidentally being assigned hyperparameters.

\subsection{Training Data Module}

\subsubsection{Every dataset covers at least one material system}
\begin{align*}
  \TrainingDatasetC &\sqsub \existsR{coversMaterial}{MaterialSystem}
\end{align*}
\emph{Example.} A Ti-Al training dataset covers the Ti-Al material system.
A dataset without any $\coversMaterialR$ link is incomplete.

\subsubsection{Every dataset covers at least one physical property}
\begin{align*}
  \TrainingDatasetC &\sqsub \existsR{coversProperty}{CoveredProperty}
\end{align*}
\emph{Example.} The Ti-Al dataset covers energies, forces, and stresses.
A dataset must declare at least one covered property.

\subsubsection{Every dataset has a provenance classification}
\begin{align*}
  \TrainingDatasetC &\sqsub \existsR{datasetProvenance}{DatasetProvenance}
\end{align*}
\emph{Example.} The Ti-Al dataset is classified as $\dlConcept{Published}$.
Every dataset must indicate whether it is published, in-house, or
augmented.

\subsubsection{Every dataset contains at least one atomic configuration}
\begin{align*}
  \TrainingDatasetC &\sqsub \existsR{hasConfiguration}{AtomicConfiguration}
\end{align*}
\emph{Example.} The Ti-Al dataset contains 15{,}000 atomic configurations.
A dataset without any configurations is empty and invalid.

\subsubsection{Every DFT calculation has settings}
\begin{align*}
  \DFTCalculationC &\sqsub \existsR{hasDFTSettings}{DFTSettings}
\end{align*}
\emph{Example.} A DFT calculation specifies PBE as the
exchange-correlation functional, a 4$\times$4$\times$4 k-point mesh, and
PAW pseudopotentials. A $\DFTCalculationC$ without $\hasDFTSettingsR$ is
incomplete.

\subsection{Benchmark Module}

\subsubsection{Every study contains at least one result}
\begin{align*}
  \BenchmarkStudyC &\sqsub \existsR{hasResult}{BenchmarkResult}
\end{align*}
\emph{Example.} A study by Smith et al.\ (2025) reports evaluation results
for MACE on Ti-Al. A $\BenchmarkStudyC$ without any $\hasResultR$ link
contains no data and is invalid.

\subsubsection{Every result evaluates exactly one trained model}
\begin{align*}
  \BenchmarkResultC &\sqsub \exactR{1}{evaluatesModel}{TrainedModel}
\end{align*}
\emph{Example.} A benchmark result evaluates the ``MACE-TiAl-v1'' model.
Each result is tied to exactly one model; comparing two models requires two
separate $\BenchmarkResultC$ instances.

\subsubsection{Every result targets at least one material system}
\begin{align*}
  \BenchmarkResultC &\sqsub \existsR{targetMaterial}{MaterialSystem}
\end{align*}
\emph{Example.} The benchmark result targets the Ti-Al material system for
testing. A result must specify what material was used for evaluation.

\subsubsection{Every result reports at least one accuracy metric}
\begin{align*}
  \BenchmarkResultC &\sqsub \existsR{hasAccuracyMetric}{AccuracyMetric}
\end{align*}
\emph{Example.} The result reports RMSE of energy (1.2~meV/atom). A
benchmark result without any accuracy metric carries no evaluation data.

\subsubsection{Every accuracy metric has exactly one type}
\begin{align*}
  \AccuracyMetricC &\sqsub \exactR{1}{metricType}{MetricType}
\end{align*}
\emph{Example.} The metric is of type RMSE. A metric cannot be both RMSE
and MAE simultaneously, nor can it lack a type.

\subsubsection{Every accuracy metric measures exactly one property}
\begin{align*}
  \AccuracyMetricC &\sqsub \exactR{1}{metricProperty}{MetricProperty}
\end{align*}
\emph{Example.} The RMSE metric measures energy. A single
$\AccuracyMetricC$ instance measures one property; separate metrics are
created for energy, force, and stress.

\subsubsection{Every result belongs to some study}
\begin{align*}
  \BenchmarkResultC &\sqsub \existsR{hasResult^{-}}{BenchmarkStudy}
\end{align*}
\emph{Example.} The benchmark result for MACE on Ti-Al belongs to the
Smith et al.\ (2025) study. A result cannot exist without a parent study.

\subsubsection{Every accuracy metric belongs to some result}
\begin{align*}
  \AccuracyMetricC &\sqsub \existsR{hasAccuracyMetric^{-}}{BenchmarkResult}
\end{align*}
\emph{Example.} The RMSE of energy metric belongs to a specific benchmark
result. An orphan metric with no parent result is invalid.

\subsection{Axioms Under Discussion}

The following axioms are candidates for inclusion but require further
discussion with domain experts.

\subsubsection{Every DFT settings belongs to some calculation}
\begin{align*}
  \DFTSettingsC &\sqsub \existsR{hasDFTSettings^{-}}{DFTCalculation}
\end{align*}
\emph{Example.} DFT settings (PBE, 4$\times$4$\times$4 k-points) belong
to a specific calculation. \emph{Discussion:} Settings could potentially be
shared across calculations as templates, in which case this axiom would be
too restrictive.

\subsubsection{Every atomic configuration belongs to some dataset}
\begin{align*}
  \AtomicConfigurationC &\sqsub \existsR{hasConfiguration^{-}}{TrainingDataset}
\end{align*}
\emph{Example.} An atomic configuration of Ti-Al belongs to the Ti-Al
training dataset. \emph{Discussion:} Configurations could belong to
multiple datasets (e.g., an augmented dataset reuses configurations from a
published one), which would still satisfy this axiom but complicates
provenance tracking.

\subsubsection{Every hyperparameter setting references a hyperparameter}
\begin{align*}
  \HyperparameterSettingC &\sqsub \existsR{forHyperparameter}{Hyperparameter}
\end{align*}
\emph{Example.} A setting ``cutoff radius = 5.0~\AA{}'' references the
hyperparameter ``cutoff radius''. \emph{Discussion:} This axiom is
straightforward but raises the question of whether settings should also be
existentially dependent on a $\TrainingRunC$ or $\TrainedModelC$.
